\section{Introducción}

El desarrollo de software moderno ha evolucionado hacia un ecosistema donde convergen la arquitectura modular, los patrones de diseño, las metodologías ágiles y la gestión eficiente de bases de datos, conformando un modelo integral de ingeniería orientado a la calidad, la escalabilidad y la reutilización. Esta convergencia responde a la necesidad creciente de las organizaciones por implementar soluciones robustas, adaptables y sostenibles, en un entorno donde la innovación tecnológica se entrelaza con la gestión estratégica de los proyectos.

En la actualidad, las arquitecturas de software no se conciben como estructuras estáticas, sino como ecosistemas dinámicos que deben responder a la evolución constante de las necesidades del negocio. Microsoft, a través de sus estrategias arquitectónicas, ha promovido estilos y patrones que buscan la escalabilidad horizontal, la independencia de componentes y la facilidad de mantenimiento. Estas estrategias se articulan en torno a principios como la separación de responsabilidades, la reutilización de servicios y la implementación de microservicios, los cuales permiten que un sistema crezca orgánicamente sin comprometer su estabilidad ni su rendimiento.

La integración de metodologías como Scrum, PMBOK y CMMI-DEV con patrones arquitectónicos reutilizables, frameworks potentes como Laravel y React, y prácticas de aseguramiento de calidad, ha permitido redefinir los estándares del desarrollo web y móvil contemporáneo. Este trabajo presenta un análisis de cómo esta sinergia metodológica y tecnológica se implementa en sistemas web empresariales, particularmente en plataformas de gestión de tickets y mesas de ayuda.

Nuestras contribuciones incluyen: (1) un análisis de la efectividad de la integración metodológica DAC con PMBOK y CMMI-DEV, (2) una evaluación de la implementación del patrón MVC en entornos Laravel-React, (3) un estudio de las prácticas de calidad del software en sistemas web empresariales, y (4) una guía de lecciones aprendidas para la implementación de arquitecturas reutilizables y modulares.